\documentclass[11pt,a4paper]{article}
\usepackage[utf8]{inputenc}
\usepackage[T1]{fontenc}
\usepackage[margin=2.5cm]{geometry}
\usepackage{enumitem}
\usepackage{parskip}
\usepackage[hidelinks]{hyperref}

\setlist{nosep,leftmargin=1.4em,itemsep=0.25em}

\title{\vspace{-1.5em}AEGIS outreach: who and when}
\author{Robin Wydaeghe \quad (Draft, May 2026)}
\date{}

\begin{document}
\maketitle
\vspace{-2em}

\section{Four purposes}
\begin{itemize}
  \item P1. Letters before 3 August. Target: 3 or 4 signed.
  \item P2. Standards allies for the 63195-2 Edition 2 window. 3 to 6 months.
  \item P3. Matched-validation co-authorship. 6 to 12 months. Replaces the ZMT track.
  \item P4. Long-term users. 1 to 3 years. The spin-off arc.
\end{itemize}

A separate lane runs in parallel: post-patent industry (Nokia, device OEMs), patient by design, not for August.

\section{P1, letters before 3 August}
A mix of academic adoption-signal letters and letters from friendly regulators, test labs, and operators. None of these need an NDA.
\begin{itemize}
  \item BIPT (Belgium). National telecom regulator. You know the EMF engineer. Credentialing letter for the file.
  \item Shanshan Wang (C2M, T\'el\'ecom Paris), Joe Wiart in cc. Wiart chairs CENELEC TC106X. GOLIAT-adjacent. Email goes to Shanshan, Wiart copied.
  \item Marta Parazzini (CNR-IEIIT Milan). Stochastic and ML dosimetry. GOLIAT. \#1 regime-matched bench group, recent mmWave \emph{Sensors} 2025 paper.
  \item Sapienza Rome (Liberti, Apollonio, or Paffi). DIET group. \#1 BioEM bench output 2023 to 2025 (19 abstracts).
  \item Gernot Schmid (Seibersdorf, AT). Address Schmid, mention Rene Hirtl by name (Hirtl led their 2023 FR2 body+hand SAR work, which is the AEGIS regime; Schmid's own recent output is ELF/LF, EV charging). IEC TC106 + IEEE ICES seats, BfS/femu/Charit\'e subcontractor pipeline, ISO 17025. Letter only, not P3.
  \item Yang Hao (Queen Mary). Body-centric mmWave. EuCAP route.
  \item Mats Gustafsson (Lund EIT/LTH). Antenna bounds. Verified email.
  \item Verkotan (Oulu, FI). Address Kai Niskala. He is the named Finnish co-convener of the IEC/IEEE 63195 standards work, in the same room as Walid El Hajj (Chair Part 3) and Teruo Onishi. FR2 OTA up to 53\,GHz, Anritsu MT8000A launch partner. CETECOM (Essen, DE) as brand-fallback runner-up, no visible 63195 footprint.
  \item Orange Innovation (Dinh-Thuy Phan-Huy). EMF-aware beam selection, RIS research. Via Margot Deruyck, 2nd-degree.
  \item Proximus, Telenet (BE operators). Brussels EMF limits much stricter than ICNIRP, so per-site compliance is harder here than elsewhere. Letter via Luc's standing relationship.
  \item ARPANSA (Australia). Only if Carolina has a real tie through her grant.
  \item GOLIAT consortium internal letter. Free, I am inside.
  \item Key stakeholders named in the IDF. List to extract.
\end{itemize}

\section{P2, the 63195-2 Ed.\ 2 window}
Process-first conversations, not letters. Your SC6 Advisory Board seat is the asset.
\begin{itemize}
  \item Akimasa Hirata (NITech). Chair ICNIRP, Chair IEEE ICES SC6. Sensitive (their group builds the obvious comparison surrogates).
  \item Kun Li (UEC Tokyo). Chairs SC6 TG6 (EPD/APD averaging).
  \item Yinliang Diao (SCAU, joint NITech). Co-Chair IEEE ICES SC6.
  \item Ilkka Laakso (Aalto). SC6 Advisory Board. Your colleague there.
  \item Teruo Onishi (NICT). Chair IEC TC106 since 2024.
  \item Walid El Hajj. Chair 63195 Part 3 (measurement side, where APD now sits in Ed.\ 2).
  \item Dragan Poljak (Split). SC6 Advisory Board. Defensive co-authorship (see P3).
\end{itemize}

\section{P3, matched validation}
The deliverable is one or two joint papers with named numbers. Without them, P2 lands as a claim.
\begin{itemize}
  \item Ericsson Research (T\"ornevik, Colombi, Xu). First call. Bern in-situ data, joint paper. Route: you to Luc.
  \item IETR Rennes (Zhadobov, Sacco). mmWave APD at 26 and 60\,GHz. Also defuses the Split overlap.
  \item CNR-IEIIT Milan (Parazzini, Bonato, Fiocchi). Their 2025 \emph{Sensors} paper on shape + skin layering is the FDTD version of AEGIS.
  \item NPL and UKHSA (Chitnis, Calder\'on, Peyman). UK regulator-affiliated metrology side. Verify 2025 \emph{Bioelectromagnetics} traceable-APD attribution before citing: agent flagged it may be IT'IS-authored, not NPL.
  \item NITech (Hirata, Kodera). Sensitive. Kodera 2024 is the empirical observation; AEGIS gives the closed-form mechanism.
  \item INRIM Turin (Bottauscio, Zilberti). Uncertainty quantification on voxel models.
\end{itemize}

\section{P4, long-term users}
Labs that would use AEGIS day to day. No letter asks here; the goal is joint work.
\begin{itemize}
  \item Shanshan Wang (C2M), Wiart cc. Stochastic surrogates need a fast forward model.
  \item INSA Lyon, Inria CITI (Mallik, Villemaud). ExposNet deep-learning EMF maps.
  \item Surrey 5G/6GIC (Elzanaty, Tafazolli). RIS for EMF reduction.
  \item Tor Vergata (Chiaraviglio, Bartoletti). RESTART 5G testbed.
  \item Thomas Eibert (TUM). Near-field rigor.
  \item Emil Bj\"ornson (KTH). Coherent-MIMO operator framing.
  \item Marco Di Renzo (Paris-Saclay). RIS. Dual with Orange Innovation. Via Margot Deruyck.
  \item CEA-LETI (Grenoble). French national RF systems lab.
  \item Oulu 6G Flagship. Testbed group that needs exposure-assessment tooling for its own experiments.
  \item ACEBR and ARPANSA (Karipidis, Loughran, Croft, Wood). ICNIRP Vice-Chair. Possibly via Carolina.
\end{itemize}

\section{Post-patent industry (a separate, patient lane)}
Not for August. Open after patent filing ($\sim$1 to 15 July).
\begin{itemize}
  \item Nokia. RAN-vendor compliance. Via Jafar Keshvari (IEEE ICES Chair), 2nd-degree through Luc and Mike Wood.
  \item Device side: Qualcomm, Apple, Samsung Mobile, MediaTek. Handset beam steering above 6\,GHz is the main APD case. Route via your IEEE ICES contacts, not cold. Qualcomm needs disclosure-timing care: the Hochwald exposure-aware-beam-selection patent (US11940477) sits there.
\end{itemize}

\section{60 days}
Weeks 1 to 2.
\begin{itemize}
  \item I send the P1 letter ask to Shanshan Wang (Wiart cc), Parazzini, Sapienza (Liberti or Apollonio), Schmid (mention Hirtl), Yang Hao, Gustafsson, Niskala (Verkotan), and to the CETECOM and Orange contacts.
  \item You send the BIPT ask and the Proximus/Telenet ask.
  \item I request the GOLIAT internal letter.
  \item You open Ericsson through Luc (matched validation, not a sales call).
  \item I ask Filip about academic letters vs commercial LOIs for the IOF jury.
\end{itemize}

Weeks 3 to 6.
\begin{itemize}
  \item I produce the matched AEGIS-vs-Sim4Life validation on a 63195-2 reference scenario (GOLIAT compute).
  \item You send a process question on 63195-2 Ed.\ 2 (Hirata or Kun Li).
  \item I email Zhadobov and Sacco about a 26 or 60\,GHz matched validation.
\end{itemize}

Weeks 7 to 8.
\begin{itemize}
  \item I open the C2M long-term conversation with Wiart.
  \item You decide on the Split approach (I lean defensive co-authorship via Poljak).
  \item Patent filing 1 to 15 July. Post-patent industry lane opens.
\end{itemize}

\section{Questions for you}
\begin{itemize}
  \item Green light on the Ericsson approach through Luc.
  \item Corrections to the P1 list, especially BIPT, the test labs, Orange, and the Belgian operators.
  \item Read on the Split approach.
  \item For Filip: do academic adoption letters carry weight, vs commercial LOIs?
  \item For Alessandro: disclosure scope per lane. My read: teaser plus password-protected demo for academic and friendly-regulator letters, NDA-only deck for industry conversations.
\end{itemize}

\end{document}
